\documentclass[a4paper, oneside]{article}
\usepackage{graphicx}
\usepackage{amsthm}
\usepackage{amsmath}
\usepackage{amssymb}
\usepackage[a4paper,
            bindingoffset=0.2in,
            left=2cm,
            right=2cm,
            top=2cm,
            bottom=2cm,
            footskip=.25in]{geometry}
\usepackage[italian]{babel}
\usepackage{pgfplots}
\usepackage{tabularx}
\usepackage{tikz}
\usepackage{wrapfig}
\usepackage{color}
\usepackage[d]{esvect}
\definecolor{page}{rgb}{0.129,0.157,0.212}
\pagecolor{page}
\color{white}
\graphicspath{ {./images/} }
\usetikzlibrary{shapes.geometric}
\usetikzlibrary{datavisualization}
\usetikzlibrary{datavisualization.formats.functions}
\usetikzlibrary{patterns}
\pgfplotsset{width=10cm,compat=1.18}

\title{Appunti di Fisica II}
\author{Tommaso Miliani}
\date{26-03-26}

\begin{document}
\newtheoremstyle{theoremEnv}
                {}          % Space above
                {}          % Space below
                {\slshape}  % Body font
                {}          % Indent amount
                {\bfseries} % Head font
                {.}         % Punctuation after head
                {\newline}  % Space after theorem head
                {}          % Theorem head spec
\theoremstyle{theoremEnv}

\newtheorem{definition}{Definizione}[section]
\newtheorem{theorem}{Teorema}[section]
\newtheorem{lemma}{Proposizione}[section]
\newtheorem{observation}{Osservazione}[section]
\newtheorem{corollary}{Corollario}[theorem]
\newtheorem{example}{Esempio}[section]
\newtheorem{remark}{Enunciato}[section]

\maketitle

\section{Prima legge di MAxwell}
\subsection{Divergenza di $\vec{B}$ }
Si diverge rispetto alla variabile del campo magnetico, ossia 
rispetto alla variabile $r$:
\begin{gather*}
    \vv{\nabla} \cdot \vv{B} = \frac{\mu_0}{4\pi}  \int_{\tau}^{} \vv{\nabla} \cdot \left(\vv{J} \times \vv{a}  \right) \ d\tau \qquad \vv{a} = \frac{(\vv{r} - \vv{r}'  )}{\left| r - r' \right|^{3} }  
\end{gather*}
Per semplicità si è eseguita una sostituzione con $\vv{a}$. Adesso si esegue la divergenza
del prodototto vettoriale
\begin{gather*}
    \vv{\nabla} \cdot (\vv{J} \times \vv{a}  ) = \frac{\partial }{\partial x_i} \varepsilon_{i, j, k} J_ja_k  = \varepsilon_{i, j, k} \frac{\partial }{\partial x_i}  J_j a_K
\end{gather*} 
Siccome la derivata è eseguita rispetto alla variabile del campo, si può portare fuori il vettore
$J$ in quanto la sua derivata è costante poiché dipende dalla variabile $r'$.
\begin{gather*}
    \vv{\nabla} \cdot  (\vv{J} \times \vv{a}  ) = J_j \varepsilon_{i, j, k} \frac{\partial }{\partial x_i} a_k  
\end{gather*}
Che assomiglia alla componente $j$ del rotore di $a$. Per trovare il rotore vero e proprio 
si scambiano $i$ e $j$ nel tensore di Ricci, così facendo si cambia di segno (che vale qualunque 
siano gli indici del tensore di Ricci), ottenendo infine
\begin{gather*}
    = J_j\left(-\varepsilon_{j, i, k} \frac{\partial }{\partial x_i} a_{k} \right) = -\vv{J} \cdot \vv{\nabla} \times \vv{a}   
\end{gather*}
Ossia il prodotto scalare tra $J$ ed il rotore di $a$. Dato che
\begin{gather*}
    \vv{a} = -\vv{\nabla} \cdot \frac{1}{\left| \vv{r} - \vv{r}'   \right| }  = + \vv{\nabla}' \cdot \frac{1}{\left| \vv{r} -\vv{r}'   \right| } 
\end{gather*} 
Allora il rotore è nullo. La legge di Biot-Savart implica che $\vv{B}$ è solenoidale. Si è dunque dimostrata la prima 
legge di Maxwell:
\begin{gather*}
    \vv{\nabla} \cdot \vv{B} = 0 \ \Longrightarrow \ \vv{\nabla}\cdot \vv{E} = \frac{\rho}{\epsilon_0}    
\end{gather*} 
Questa equazione mi die che non esistono densità di carica magnetiche (quindi non 
esistono monopoli magnetici). Il campo magnetico definito fino ad ora 
non è generato da qualcosa di analogo alle cariche elettriche. Segue dunque che 
\begin{gather*}
    \Phi _S(\vv{B} ) = 0
\end{gather*}
qualunque sia la superficie $S$ considerata in qualunque punto dell'universo e 
qualunque sia $\vv{B}$ generato. Le linee di campo di $\vv{B}$ sono chiuse: non esiste 
allora una sorgente (di cariche) per $\vv{B}$, come invece accade per le cariche elettriche.   
Il concetto di sorgenti nei campi si estende alla quantità che genera il campo. $\vv{B}$ non ha monopoli
ma chiaramente le correnti elettriche sono sorgenti per $\vv{B}$. Ci si aspetta dunque 
che nell'altra equazione di maxwell compaia il vettore $\vv{J}$ e ci si aspetta dunque 
che non sia un campo irrotazionale.    
\subsection{Rotore di $\vec{B}$ }
Il rotore 
\begin{gather*}
    \vv{\nabla}\times \vv{B} = \int_{\tau} \vv{\nabla} \times (\vv{J} \times \vv{a}  ) \ d\tau    
\end{gather*}
Dove $\vv{a}$ è come definito prima. Si fissa una coordinata alla volta per il conto:
\begin{gather*}
    \left\{ \vv{\nabla} \times \left(\vv{J} \times \vv{a}  \right) \right\}_1 = \varepsilon_{1, j, k} \frac{\partial }{\partial x_j} (\vv{J} \times \vv{a}  )_k = \varepsilon_{1, j, k} \frac{\partial }{\partial x_j} \varepsilon_{k, l, m} J_l a_m  
\end{gather*}
Si esplicitano ora gli indici non nulli: 
\begin{align*}
    \left\{ \vv{\nabla} \times \left(\vv{J} \times \vv{a}  \right) \right\}_1 &= \varepsilon_{1, 2, 3} \frac{\partial }{\partial x_2} \left(\varepsilon_{3, l , m} J_l a_m\right)  - \varepsilon_{1, 3, 2} \frac{\partial }{\partial x_3} \varepsilon_{2 ,l , m} J_l a_m    \\
    &= \frac{\partial }{\partial x_2} \left(\varepsilon_{3, l , m} J_l a_m\right)  - \frac{\partial }{\partial x_3} \varepsilon_{2 ,l , m} J_l a_m     \\
    &= \varepsilon_{3, 2, 1} \frac{\partial }{\partial x_2} J_2 a_1 + \varepsilon_{3, 1, 2} \frac{\partial }{\partial x_2} J_1 a_2 - \varepsilon_{2, 3, 1} \frac{\partial }{\partial x_3} J_3 a_1 - \varepsilon_{2, 1, 3} \frac{\partial }{\partial x_3} J_1 a_3  \\
    &= -\frac{\partial }{\partial x_2} J_2 a_1 + \frac{\partial }{\partial x_2} j_1 a_2 - \frac{\partial }{\partial x_3} J_3 a_1 + \frac{\partial }{\partial x_3} J_1 a_3 \\
    &= J_1\left(\frac{\partial }{\partial x_2} a_2 + \frac{\partial }{\partial x_3} a_3  \right) - J_2 \frac{\partial }{\partial x_2} - J_3 \frac{\partial }{\partial x_3} a_1
\end{align*}
Dove l'ultimo passaggio vale perché $J$ non è funzione di $r$. Si somma e sottrae adesso un termine: 
\begin{align*}
    \left\{ \vv{\nabla} \times \left(\vv{J} \times \vv{a}  \right) \right\}_1     &= J_1\left(\frac{\partial }{\partial x_2} a_2 + \frac{\partial }{\partial x_3} a_3  \right) - J_2 \frac{\partial }{\partial x_2} - J_3 \frac{\partial }{\partial x_3} a_1 + J_1 \frac{\partial }{\partial x_1} a_1 -  J_1 \frac{\partial }{\partial x_1}a_1 \\
    &=   J_1 \left(\frac{\partial }{\partial x_1} a_1 + \frac{\partial }{\partial x_2} a_2 + \frac{\partial }{\partial x_3} a_3   \right) - \left(J_1 \frac{\partial }{\partial x_1} a_1 + J_2 \frac{\partial }{\partial x_2}a_1 + J_3 \frac{\partial }{\partial x_3} a_1   \right)
\end{align*}
Il primo termine è la divergenza di $\vv{a}$, mentre il secondo termine (che si chiama $\mathcal{C}$) necessita ancora di 
alcuni passaggi.  Si cambia segno al termine e si deriva secondo la seguente espressione (ossia la derivata
del prodotto rispetto a $x'$):
\begin{gather*}
    J_1 \frac{\partial }{\partial x_1'} a_1 = \frac{\partial }{\partial x_1'}(J_1 a_1) - a_1 \frac{\partial }{\partial x_1'}J_1     
\end{gather*}
Dunque, se si scambiano le parziali con $x$ con le parziali per $x'$ si ha 
\begin{gather*}
    \mathcal{C}  = \frac{\partial }{\partial x_1'}(J_1 a_1) -a_1 \frac{\partial }{\partial x_1'}J_1 + \frac{\partial }{\partial x_2'}(J_2 a_1) - a_1 \frac{\partial }{\partial x_2'}J_2 + \frac{\partial }{\partial x_3'}(J_3 a_1) - a_1\frac{\partial }{\partial x_3'} J_3       
\end{gather*}
Si può dunque riscrivere
\begin{gather*}
    \mathcal{C} = \vv{\nabla}' \cdot \left(a_1 \vv{J} \right) - a_1 \vv{\nabla}' \cdot \vv{J}  
\end{gather*}
L'ultimo termine è uguale a zero in statica (poiché il campo deve essere solenoidale). 
Si è dunque trovato che
\begin{gather*}
    \left\{ \vv{\nabla} \times \left(\vv{J} \times \vv{a}  \right) \right\}_1 = J_1 \vv{\nabla} \cdot \vv{a} + \vv{\nabla}' \cdot (a_1 \vv{J} )   
\end{gather*}
Si può adesso calcolare la componente:
\begin{gather*}
    (\vv{\nabla} \times \vv{B})_1  = \frac{\mu_0}{4\pi} \int_{\tau} J_1 \vv{\nabla} \cdot \vv{a} \ d\tau + \frac{\mu_0}{4\pi} \int_{\tau} \vv{\nabla}' \cdot (a_1 \vv{J} )  \ d\tau  = \frac{\mu_0}{4\pi} \int_{\tau} J_1 \vv{\nabla} \cdot \vv{a} \ d\tau +\frac{\mu_0}{4\pi} \int_{S} \hat{n} \cdot (a_1 \vv{J} ) \ dS  
\end{gather*}
Dato che nessun punto della superficie tocca il circuito, $\vv{J}$ è identicamente nulla su tutti i punti 
della superficie $S$ definita perché in statica $\vv{J}$ è tangenziale in tutti i punti della superficie del filo 
in quanto se fosse presente una componente normale allora ci sarebbe accumulo di carica 
sulla superficie del filo, dunque il secondo pezzo è nullo. Allora 
\begin{align}
    \vv{\nabla} \times \vv{B} = \frac{\mu_0}{4\pi} \int_{\tau} \vv{\nabla} \times (\vv{J} \times \vv{a}  ) \ d\tau = \frac{\mu_0}{4\pi} \int_{\tau} \vv{J} \cdot \vv{\nabla} \cdot \vv{a} \ d\tau         
\end{align}  
Se si utilizzasse la coordinata campo:
\begin{gather*}
    = \frac{\mu_0}{4\pi} \int_{\tau} \vv{J}\left(\vv{\nabla} \cdot \left(\vv{\nabla}' \frac{1}{\left| r - r' \right| } \right) \right)  \ d\tau = - \frac{\mu_0}{4\pi} \int \vv{J}(r') \left(\nabla^{2} \frac{1}{\left| r - r' \right| }\right) \ d\tau  
\end{gather*}
Si è dunque concluso che 
\begin{gather*}
    \nabla^{2} \frac{1}{\left| r - r' \right| } = -4\pi \delta (\vv{r} - \vv{r}'  )
\end{gather*}
Nel caso in cui nel campo vettoriale $\vv{r}$ sia la posizione allora ha le dimensioni 
di 1 su volume:
\begin{gather*}
    \vv{\nabla} \times \vv{B} = - \frac{\mu_0}{4\pi} \int \vv{J}(\vv{r}' ) \left(-4\pi \delta (\vv{r} - \vv{r}'  )\right) \ d\tau = \mu_0 \int \vv{J}(\vv{r}' ) \delta (\vv{r} - \vv{r}'  ) \ d\tau      
\end{gather*}
Questo integrale dipende dal dominio di integrazione: fa zero se non contiene il punto di
divergenza della $\delta$. Se si calcola il rotore di $\vv{B}$ in un punto in 
cui esiste $\vv{J}$ vale
\begin{gather*}
    \mu_0 \vv{J}(\vv{r} ) 
\end{gather*}  
se si va a calcolare il rotore in un punto in cui non c'è corrente, allora
$\vv{r}$ non fa pare del dominio di integrazione e dunque è nullo. Si è 
trovata allora la quarta equazione di MAxwell: 
\begin{align}
    \vv{\nabla} \times \vv{B} = \mu_0 \vv{J}   
\end{align} 
Allora $\vv{B}$ è solenoidale che è generato da $\vv{J}$ che entra nelle equazioni 
di MaxWell per mezzo del rotore e $\vv{J}$ è la corrente ed è la sorgente di $\vv{B}$.

\section{Trattazione equazioni di MaxWell}
Le equazioni di Maxwell
\begin{align}
    &\vv{\nabla} \cdot \vv{E} = \frac{\rho}{\epsilon_0} \\
    &\vv{\nabla} \cdot \vv{B} = 0 \\
    &\vv{\nabla} \times \vv{E}= 0 \\
    &\vv{\nabla} \times \vv{B} = \mu_0 \vv{J}         
\end{align}
descrivono tutte quantità disgiunte tra loro (nella statica). Queste si possono 
catalogare in equazioni omogenee (seconda e terza) o in equazioni disomogenee o forzate (come 
la prima e la quarta). Da quelle omogenee si ricava una semplificazione 
della fisica: infatti la terza implica che $\vv{E} = - \vv{\nabla}V$, ossia 
una quantità vettoriale è descrivibile tramite il gradiente di una funzione scalare. 
Invece la seconda implica che $\vv{B} = \vv{\nabla} \times \vv{A}$, che si tratterà ora.    
Si vuole dimostrare ora che se
\begin{gather*}
    \vv{B} = \vv{\nabla} \times \vv{A} \ \Longrightarrow \ \vv{\nabla} \cdot \vv{B} = 0     
\end{gather*}
Si esprime allora 
\begin{gather*}
    \vv{\nabla} \cdot (\vv{\nabla} \times \vv{A}  ) = \frac{\partial }{\partial x_i} (\vv{\nabla} \times \vv{A}  )_i = \frac{\partial }{\partial x_i}\varepsilon_{i, j, k} \frac{\partial }{\partial x_j} A_k = \varepsilon_{i,j, k} \frac{\partial }{\partial x_i} \frac{\partial }{\partial x_j}A_k      
\end{gather*}
Per poter svolgere questo conto si "smezza" l'ultimo termine: 
\begin{gather*}
    \varepsilon_{i,j, k} \frac{\partial }{\partial x_i} \frac{\partial }{\partial x_j}A_k  = \frac{1}{2} \varepsilon_{i,j, k} \frac{\partial }{\partial x_i} \frac{\partial }{\partial x_j}A_k  + \frac{1}{2} \varepsilon_{i,j, k} \frac{\partial }{\partial x_i} \frac{\partial }{\partial x_j}A_k  
\end{gather*} 
Adesso nel secondo termine si scambiano gli indici $i$ e $j$ (nella parte antisimmetrica):
\begin{gather*}
    \varepsilon_{i,j, k} \frac{\partial }{\partial x_i} \frac{\partial }{\partial x_j}A_k  = \frac{1}{2} \varepsilon_{i,j, k} \frac{\partial }{\partial x_i} \frac{\partial }{\partial x_j}A_k  - \frac{1}{2} \varepsilon_{j, i k} \frac{\partial }{\partial x_i} \frac{\partial }{\partial x_j}A_k  
\end{gather*}
Sempre nel secondo termine si scambiano gli indici di derivazione. Scambiando gli indici 
nella parte antisimmetrica si moltiplica per un segno meno, mentre se si scambiano gli indici 
nella parte simmetrica si ottiene sempre un segno più.
\begin{gather*}
    \frac{1}{2} \varepsilon_{i,j, k} \frac{\partial }{\partial x_i} \frac{\partial }{\partial x_j}A_k  - \frac{1}{2} \varepsilon_{j, i k} \frac{\partial }{\partial x_j} \frac{\partial }{\partial x_i}A_k  
\end{gather*}
Si scambiano adesso $j = l$ e $i = m$ per il primo termine e $i =l $ e $j = m$ nel secondo termine:
\begin{gather*}
    \frac{1}{2} \varepsilon_{l, m, k} \frac{\partial }{\partial l}\frac{\partial }{\partial m} A_k + \frac{1}{2} \varepsilon_{l, m, k} \frac{\partial }{\partial l} \frac{\partial }{\partial m} A_k = 0    
\end{gather*}
Dunque si è dimostrata la condizione iniziale. Adesso si vuole utilizzare 
il rotore di $\vv{A}$ e legarlo al vettore che esprime la corrente 
secondo la seguente condizione:
\begin{align}
    \vv{\nabla} \times (\vv{\nabla} \times \vv{A}) = \mu_0 \vv{J}    
\end{align} 
Per dimostrarla adesso si fa, come fatto prima, una componente alla volta:
\begin{align*}
    \left(\vv{\nabla} \times (\vv{\nabla} \times \vv{A}  ) \right)_1 &= \varepsilon_{i, j, k} \frac{\partial }{\partial x_i} \left(\vv{\nabla} \times \vv{A}  \right)_k = \varepsilon_{1, j, k} \frac{\partial }{\partial x_j} \varepsilon_{k, l, m} \frac{\partial }{\partial x_l} A_m \\
    &= \varepsilon_{3, 2, 1} \frac{\partial }{\partial x_2} \frac{\partial }{\partial x_2} A_1 + \frac{\partial }{\partial x_2}       \frac{\partial }{\partial x_1} A_2 - \varepsilon_{2, 1 ,3} \frac{\partial }{\partial x_3}  \frac{\partial }{\partial x_1} A_3 - \varepsilon_{2,3,1 } \frac{\partial }{\partial x_3}  - \varepsilon_{2, 3, 1} \frac{\partial }{\partial 3} \frac{\partial }{\partial 3} A_1 \\
    &= -\frac{\partial }{\partial x_2} \frac{\partial }{\partial x_2} A_1 + \frac{\partial }{\partial x_2} \frac{\partial }{\partial x_1} + \frac{\partial }{\partial x_3} \frac{\partial }{\partial x_1} A_3 - \frac{\partial }{\partial x_3} \frac{\partial }{\partial x_3} A_1 + \frac{\partial ^{2}}{\partial x_1^{2}} A_1 - \frac{\partial ^{2}}{\partial x_1^{2}} A_1          \\
    &= \frac{\partial }{\partial x_1} \left(\frac{\partial }{\partial x_i}   A_1 \right) - \frac{\partial }{\partial x_i} \frac{\partial }{\partial x_i}A_1   = \vv{\nabla} \cdot  \left(\vv{\nabla} \cdot A \right) - \nabla^{2}\vv{A}  
\end{align*}
Una cosa importante è il $\nabla^{2}$ di un campo vettoriale, esso è definito in modo 
divereso per campi scalari e per campi vettoriali:
\begin{align*}
    \nabla^{2} f &= \frac{\partial }{\partial x_i} \frac{\partial }{\partial x_i}f = \vv{\nabla} (\vv{\nabla}f ) \\
    \nabla^{2} \vv{A} &= \nabla^{2}A_x \hat{x} + \nabla^{2} A_y \hat{y} + \nabla^{2} A_z \hat{z}   
\end{align*}
In magnetostatica si utilizza il \textbf{Gauge di Coulomb}, ossia un vincolo che si pone 
sul campo vettoriale $\vv{A}$, che dice che gli unici $\vv{A}$ possibili sono quelli solenoidali, 
per cui $\vv{\nabla} \cdot \vv{A} = 0$. Se si sceglie il Gauge di Coulomb, allora 
\begin{align}
    (\vv{\nabla} \times (\vv{\nabla} \times \vv{A}  ) ) = -\nabla^{2}\vv{A} 
\end{align}
E dunque 
\begin{align}
    \nabla^{2} \vv{A} = - \mu_0 \vv{J}  
\end{align}

\subsection{Problema di Ampere}
In elettrostatica, l'unione della prima e terza equazione di Maxwell porta 
all'equazione di Poisson:
\begin{align}
    \nabla^{2} V = - \frac{\rho}{\epsilon_0} \Longleftrightarrow V(\vv{r} ) = \frac{1}{4\pi\epsilon_0} \int_{\rho(\vv{r} )}^{\left| \vv{r} - \vv{r}'   \right| }  \ d\tau
\end{align}
Lo stesso vale per la magnetostatica  unendo la seconda e la quarta equazione di Maxwell:
\begin{align}
    A_i(\vv{r} ) = \frac{\mu_0}{4\pi} \int_{\tau} \frac{J_i(\vv{r}' )}{ \left| \vv{r} - \vv{r}'   \right| } \ d\tau 
\end{align}
Allora
\begin{align}
    \vv{A}(\vv{r} ) = \frac{\mu_0}{4\pi} \int_{\tau} \frac{\vv{J}(\vv{r}' ) }{\left| \vv{r} - \vv{r}'   \right| }  \ d\tau
\end{align}
Che è vero se e solo se $\vv{\nabla} \cdot \vv{A} = 0$. Che aiuta a risolvere i 
problemi di elettrostatica in maniera più semplice. La priam equazione di Maxwell 
esiste sia in forma integrale (Gauss) che vettoriale. Dunque questo porta
al \textbf{Problema di Ampere}, che lo si calcola tramite la forma integrale della prima equazione 
di Maxwell:
\begin{gather*}
    \int_{\tau} \vv{\nabla} \times \vv{B} \hat{u} \ dS    
\end{gather*}
Secondo una superficie qualsiasi orientata in un qualsiasi modo. Si utilizza allora 
il teorema di Stokes per cui
\begin{align}
        \int_{\tau} \vv{\nabla} \times \vv{B} \hat{u} \ dS = \oint_\gamma \vv{B} \cdot \vv{dl}  
\end{align}
secondo la regola della mano destra. Questo lo si scrive anche come 
\begin{gather*}
    \int_{S} \vv{J} \cdot \hat{u} \ dS   
\end{gather*}
Allora 
\begin{gather*}
    \int_{\tau} \vv{\nabla} \times \vv{B} \hat{u} \ dS = \oint_\gamma \vv{B} \cdot \vv{dl}  =     \int_{S} \vv{J} \cdot \hat{u} \ dS 
\end{gather*}
Questo è uguale a dire 
\begin{gather*}
    \mu_0 \int_{S} \vv{J} \cdot \hat{u} \ dS = \mu_0 I_c   
\end{gather*}
dove $I_c$ prende il nome di \textbf{corrente concatenata al circuito $\gamma$}. Se $\gamma$
si concatena al circuito esiste almeno una sezione del circuito che passa per $\gamma$. Se si aggiungono 
più superfici attraverso la quale passa $I$: 
\begin{align}
        \int_{\tau} \vv{\nabla} \times \vv{B} \hat{u} \ dS = \oint_\gamma \vv{B} \cdot \vv{dl}  =     \int_{S} \vv{J} \cdot \hat{u} \ dS = \mu \sum_i I_i = \mu \sum I_{c_i}
\end{align}
\begin{example}
    Supponendo di avere un cilindro con $\vv{J}$, in coordinate
    cilindriche il campo mangetico ha espressione:
    \begin{align}
        \vv{B}(r, \theta, z) = B_r(r, \theta, z) \hat{r}  +  B_\theta(r, \theta, z) \hat{\theta}+  B_z(r, \theta, z) \hat{z}
    \end{align} 
    Per la simmetria rimane solo il termine $\theta$ poiché sono linee che girano intorno 
    al filo (allora non dipende dalla coordinata $\theta$ né da $z$ in quanto è invariante 
    rispetto alla traslazione e per invarianza sulla rotazione):
    \begin{gather*}
        \vv{B} = B_\theta(r) \hat{\theta}  
    \end{gather*}
    Circonferenze concentriche al filo per calcolare il teorema di Ampere, 
    su tutti i punti della circonferenza si ha la stessa espressione, allora 
    \begin{gather*}
        \oint_\gamma \vv{B} \cdot \vv{dl} = 2\pi r B_\theta  
    \end{gather*}
    Il segno dipende dal versore $\theta$. Allora, se il filo avesse 
    raggio $R$:
    \begin{gather*}
        I = J\pi R^{2}
    \end{gather*}
    Il teorema di Ampere mi dice che
    \begin{gather*}
        \oint_\gamma \vv{B} \cdot \vv{dl} = \left\{\begin{array}{l}
            \displaystyle r > R \ \Longrightarrow \ 2\pi rB\theta = \mu_0 I \\
            \displaystyle r < R \ \Longrightarrow \ 2\pi rB_\theta = \mu_0\pi J r^{2} 
        \end{array}\right.  
    \end{gather*}
    Si trova allora che $B_\theta$ cresce linearmente fino a $R$ e decresce 
    come $\sim 1 /r^{2}$ dopo $R$.
\end{example}



\end{document}